
%\needspace{4\baselineskip}
\color{uniblue}{\subsection{Междуфазная токовая отсечка (МФТО)}}
\color{black}

\begin{enumerate}[label=\arabic{section}.\arabic{subsection}.\arabic{enumi}, labelsep=4pt, leftmargin=0pt, itemindent=57pt, itemsep=0pt, parsep=0pt] % если что можно крутить itemsep - ЗДЕСЬ нет подпунктов, только пункты - поэтому parsep=0pt а не parsep=5pt, как с подпунктами

\item
Для резервирования действия защит при близких КЗ в устройстве предусмотрена ненаправленная междуфазная токовая отсечка (МФТО), работающая по факту превышения уставок не менее чем двумя из трех фазных токов.

\item
Ввод в работу МФТО осуществляется программным ключом \textbf{«МФТО»}.

\item
Ток срабатывания МФТО задается с помощью уставки \textbf{«Iср МФТО»}. Время срабатывания задаётся уставкой \textbf{«Тср МФТО»}. 

\item
Для предотвращения неправильной работы МФТО предусмотрена блокировка при протекании через трансформаторы тока двух выключателей сквозного тока внешнего КЗ, вызывающего неодновременное насыщение трансформаторов. Блокировка вводится программным ключом \textbf{«БСТО МФТО»}.

\item
МФТО может быть оперативно выведена из работы от внешнего сигнала \textit{«Вывод МФТО»}, назначенного на соответствующий дискретный вход или функциональную клавишу.

\item
Функциональная схема логики МФТО приведена на рисунке \ref{pict_mfto:MFTO}.

\vspace{2mm}
\begin{figure}[H]
\centering
\includegraphics[width=1\textwidth,height=1\textheight,keepaspectratio]{\fbpath/561.0171 МФТО/img/MFTO.pdf}
\captionof{figure}{Функциональная схема МФТО}\label{pict_mfto:MFTO}
\end{figure}

\item
Уставки МФТО представлены в таблице \ref{desc_mfto:tbl1}.

\small
\begin{longtable}{|>{\centering\arraybackslash}m{5.3cm}|>{\centering\arraybackslash}m{3.3cm}|>{\centering\arraybackslash}m{4.2cm}|>{\centering\arraybackslash}m{1.8cm}|>{\centering\arraybackslash}m{1cm}|}
\caption{Параметры для настройки МФТО\hfill\vspace{-0.5\baselineskip}}\label{desc_mfto:tbl1}\\ 
\hline
\rowcolor{gray!30}
Параметр & Условное обозначение на схеме & Значение/ Диапазон & Единица измерения & Шаг \\ 
\hline
\endfirsthead
\caption*{\hspace{3pt}\emph{Продолжение таблицы \ref{desc_mfto:tbl1}\hfill\vspace{-0.5\baselineskip}}} \\ % сделано по ГОСТ 2.105 п.6.8.7
\hline
\rowcolor{gray!30}
Параметр & Условное обозначение на схеме & Значение/ Диапазон & Единица измерения & Шаг \\ 
%\hline
\endhead

%===m>МФТО|МФТО
\raggedright {'Description': 'МФТО'} & \centering Pereklyuchatelvkhodovpoustavkebool2S & \centering Введено \\ Выведено & \centering --- & \centering \arraybackslash --- \\
\hline
\raggedright {'Description': 'Iср МФТО'} & \centering Relesktomvozvratamaxfloat1xsr & \centering 0,04 ... 30,00 & \centering о.е. & \centering \arraybackslash 0,01 \\
\hline
\raggedright {'Description': ''} & \centering Pereklyuchatelvkhodovpoustavkebool1S & \centering Введено \\ Выведено & \centering --- & \centering \arraybackslash --- \\
\hline
\raggedright {'Description': 'Тср МФТО'} & \centering NezavisimoeRVnasrabatyvanie1Tvremya & \centering 0 ... 300000 & \centering мс & \centering \arraybackslash 5 \\
\hline
%===m
\end{longtable}
\normalsize

\end{enumerate}
